\subsection{Solvable closures of finitely generated fields}

\begin{claim}{F.43}{Proved}
Let $K/\Q$ be a finitely generated field, and let
$K^{\mathrm{solv}}$ be its maximal Galois prosolvable
extension in a chosen algebraic closure. Every subfield $k_0$
of a finitely generated extension of $K^{\mathrm{solv}}$
satisfies closed-point-preserving reconstruction: for every
pair of finite-type $k_0$-schemes $X,Y$, the canonical map
\[
 \operatorname{Hom}_{k_0}(X^{\mathrm{awn}},Y)
 \longrightarrow
 \pi_0\operatorname{Hom}^{\mathrm{cp}}_{BG_{k_0}}
       (X_{\mathrm{et}},Y_{\mathrm{et}})
\]
is bijective.

The class of fields in this statement is equivalently described
by the following sufficient arithmetic condition: there is a
finitely generated subfield $F_0\subset k_0$ such that
$k_0/F_0$ is algebraic and the Galois group of its normal
closure has an open prosolvable subgroup. This class strictly
contains F.37's sufficient class: $\Q(t)^{\mathrm{solv}}$
does not embed in a finitely generated extension of
$\Q^{\mathrm{solv}}$. No necessity of this sufficient
condition, and no classification as in F.9, is asserted.
\end{claim}
\begin{proof}
Write $k=K^{\mathrm{solv}}$ and
$D(E^\times)=\bigcap_{n\ge1}(E^\times)^n$.
We first prove reconstruction over every finitely generated
$L/k$, verifying the arithmetic inputs of F.33--F.37.
The final subfield descent needs a separate argument, since
$k_0k/k_0$ need not be algebraic when $K/\Q$ has positive
transcendence degree.

\proofparagraph{Constants and finite Galois composita.}
For a Galois extension $E/F$ and an extension $L/F$ in a
common overfield, restriction identifies
\[
 \operatorname{Gal}(EL/L)
       \simeq\operatorname{Gal}(E/(E\cap L)).                 \tag{F.43.1}
\]
This also holds for infinite Galois extensions. The
restriction image is compact and hence closed, and its
fixed field in $E$ is $E\cap L$; the Galois correspondence
gives the assertion. This is precisely
\citet[Proposition~7.15 and Corollary~7.16, p.~99]{F_MilneFT2022}.
No normality of $L/F$ is required.

Let $C=K\cap\overline{\Q}$. This is a number field, by
finiteness of the relative algebraic closure in a finitely
generated extension \citep[Lemma~9.26.11]{F_StacksFields}.
The extension $K/C$ is regular. Explicitly, a finite Galois
$N/C$ in $\overline{\Q}$ has $N\cap K=C$, so (F.43.1)
makes $N$ and $K$ linearly disjoint over $C$. We claim
\[
              k\cap\overline{\Q}=C^{\mathrm{solv}}.           \tag{F.43.2}
\]
Every finite Galois solvable $N/C$ gives a Galois solvable
$KN/K$, proving one inclusion. Conversely, take
$a\in k\cap\overline{\Q}$ and a finite Galois solvable
$E/K$ in $k$ containing $a$. Let $N/C$ be the normal
closure of $C(a)/C$. Regularity makes all $C$-conjugates
of $a$ its $K$-conjugates. Normality of $E/K$ therefore
puts $N$ inside $E$. The group
$\operatorname{Gal}(N/C)\simeq\operatorname{Gal}(KN/K)$
is a quotient of the solvable group
$\operatorname{Gal}(E/K)$. Hence $a\in C^{\mathrm{solv}}$,
proving (F.43.2).

For arbitrarily large $n\ge5$, Hilbert specialization over
the number field $C$ gives an $S_n$-Galois extension $N/C$,
by the Hilbert inputs already used in F.18 and F.34
\citep[Corollary~3.2.2, Proposition~3.3.1 and
Theorem~3.4.1]{F_Serre2007}. Its intersection with
$C^{\mathrm{solv}}$ is $N^{A_n}$: every solvable quotient
of $S_n$ kills $A_n$, and the quadratic quotient already
lies in the solvable closure. Thus the residual degrees
$[NC^{\mathrm{solv}}:C^{\mathrm{solv}}]=|A_n|$ are
unbounded. Equations (F.43.1)--(F.43.2) imply
\[
 [k\overline{\Q}:k]
       =[\overline{\Q}:C^{\mathrm{solv}}]=\infty.             \tag{F.43.3}
\]
If $L/k$ is finitely generated, its relative algebraic
subfield over $k$ is finite, again by
\citet[Lemma~9.26.11]{F_StacksFields}. Consequently $L$
cannot contain $\overline{\Q}$. In particular, there
exists a finite Galois \emph{number field} $N_0/\Q$
not contained in $L$. Supplying a missing number field,
and not just a finite geometric extension of $K$, is
essential to the arithmetic-place construction below.

\proofparagraph{Divisibility and fixed finite descent fields.}
The field $k$ contains $\mu_\infty$ and $k^\times$ is
divisible. To see the latter, put $a\in k$ in a finite
Galois solvable $E/K$, adjoin the required roots of unity,
and adjoin $n$-th roots of every $K$-conjugate of $a$.
The radical group over the field containing the roots
of unity is abelian; the resulting normal extension
over $K$ is still solvable. It lies in $k$.

For every finitely generated $F/\Q$,
$D(F^\times)=\{1\}$. On a normal projective model over
the number field of algebraic constants, the divisor
of an infinitely divisible element is divisible by
every positive integer and hence is zero. The element
and each of its roots are constants. Their finite
valuations then vanish, and the number-field unit
theorem reduces the claim to the intersection of all
power subgroups of a finitely generated abelian group.
That intersection is trivial.

For a finite extension $E/k$, enlarge it to $Mk/k$ with
$M/K$ finite Galois. Put $B=M\cap k$. The extensions
$M/B$ and $k/B$ are Galois and disjoint, so their
restrictions commute in the compositum. For a fixed
$a\in Mk$, choose a finite $S/K$ inside $k$ containing
$B$ and all scalar coefficients needed to write $a$,
and put $F=MS$. This field is finitely generated over
$\Q$ and is fixed independently of any radical level.
If $z^n=a$ with $z\in Mk$, then, for
$\sigma\in\operatorname{Gal}(Mk/k)$ and
$\tau\in\operatorname{Gal}(Mk/F)$,
\[
 \tau z/z\in\mu_n\subset k,\qquad
 \tau(\sigma z/z)
   =\frac{\sigma(\tau z/z)}{\tau z/z}\,\sigma z/z
   =\sigma z/z.
\]
Thus $\sigma z/z\in F$, and
$\sigma a/a\in(F^\times)^n$. If $a\in D((Mk)^\times)$,
this holds for every $n$, so $D(F^\times)=1$ gives
$\sigma a=a$. Hence $a\in k$. The converse follows
from divisibility of $k^\times$, and therefore
\[
                   D(E^\times)=k^\times
                   \quad(E/k\text{ finite}).                \tag{F.43.4}
\]
For finitely generated $L/k$, a normal projective
$k$-model has as its global functions the finite
relative algebraic closure $k'$ of $k$ in $L$.
The divisor argument places an infinitely divisible
element and all its roots in $k'$. Applying (F.43.4)
there proves
\[
                  D(L^\times)=k^\times.                    \tag{F.43.5}
\]
This holds after every finite extension of $L$ as well.

Given any finite collection of data over $L$, choose
a finitely generated $F/\Q$ inside $L$ containing
$K$, generators for $L/k$, and all coefficients and
relations in that collection. Then
\[
                  K\subset F,\qquad L=Fk,
       \qquad L/F\text{ is Galois prosolvable}.              \tag{F.43.6}
\]
Indeed the last extension is the base change of $k/K$.
Finite extensions of $L$ and finitely many elements
in them can be written in $Mk$ for a finite $M/F$,
after enlarging $F$ by finitely many elements of $k$.
Those elements are algebraic over the finitely generated
field $K$, so $F$ and $M$ remain finitely generated
over $\Q$. This verifies the finite-generation
predicate at every descent below. In particular no
claim that $k/\Q$ is normal is being made.

\proofparagraph{Finite arithmetic quotient detection.}
Descend a smooth geometrically connected curve $U/L$
and a connected finite $G$-torsor, including the whole
action and the torsor identity, to an $F$ as in
(F.43.6). If $M/F$ is the finite Galois constant
field of that cover, connectedness after base change
gives $M\cap L=F$. Write
$1\to N\to G\to P\to1$, with
$P=\operatorname{Gal}(M/F)$.
The geometric pencil construction of F.24, at unbounded
degrees $n\ge5$, has Galois closure with group
\[
 H\subseteq(G\times_P\cdots\times_PG)\rtimes S_n.
\]
Its ordered source uses one common constant field $M$.
The spanning-tree gauge and the commutators supported
at at most two coordinates in F.24 show that
$A_n\subset H^\infty$ and that
$H^\infty\cap N^n\to N$ is onto at each coordinate.
These are geometric calculations over a finitely
generated characteristic-zero field, which embeds
in $\C$. Thus the same branch-loop and comparison
proof applies; no full wreath-product assertion is
required.

The field $F/\Q$ is Hilbertian, by the rational-function
and finite-extension Hilbert theorems used in F.34.
Specializing with the full group $H$ gives
$\Omega/F$. By (F.43.6),
\[
       J=\operatorname{Gal}(\Omega L/L)\supset H^\infty,
       \qquad J\longrightarrow P\text{ is onto}.             \tag{F.43.7}
\]
The first inclusion holds because $\Omega\cap L/F$
is solvable; the second follows from $M\cap L=F$.
A coordinate stabilizer in $J$ surjects onto $G$.
The support-two commutators supply $N$ there, and
a literal even permutation moves the selected block
back without changing a prescribed $P$-image.
This proves detection of any prescribed finite
\emph{arithmetic} quotient at a closed point of $U$.
The specialized degree-$n$ fibre is a single closed
point because its permutation image contains the
transitive $A_n$. In particular such closed points
have unbounded degrees.

For a smooth quasiprojective source of dimension
at least two, apply the geometric-component-preserving
Bertini reduction of F.29. The quasi-finite map of
each geometric component of the cover has equidimensional
nonempty fibres, so the exact hypothesis of
\citet[Theorem~1.5]{F_PoonenSlavov2020} holds.
Successive good hyperplanes preserve every geometric
component and its arithmetic permutation. Their good
open subsets have $L$-points since $L$ is infinite.
The resulting curve retains the entire arithmetic
finite quotient, and the preceding detection applies.

\proofparagraph{Punctured-line integrality and the normality check.}
Let $U\subset\mathbb G_{m,L}$ be nonempty open, with
irreducible puncture polynomials
$P_0=T,P_1,\ldots,P_r$. Suppose a compatible Kummer
class has actual multiplicative values at every point
over every finite extension of $L$. Choose
$x_0\in U(L)$ with actual value $b_0$. The vanishing
Picard group of this localized principal ideal
domain gives the expression
\[
 \kappa(b_0)+
 \sum_j\lambda_j\bigl([P_j(T)]-[P_j(x_0)]\bigr),
 \qquad\lambda_j\in\widehat{\Z}.                              \tag{F.43.8}
\]
Choose a finitely generated $F_0/\Q$ containing
$K$, the displayed coefficients and normalization,
and with $L=F_0k$. For a transcendence basis
$t_1,\ldots,t_d$ of $F_0/\Q$, take a finite Galois
$M/\Q(t_1,\ldots,t_d)$ containing $F_0$, a missing
number field $N_0$ from (F.43.3), and all the
$P_j$-roots. The finite extension $Mk/L$ is Galois.
There is $\sigma\in\operatorname{Gal}(Mk/L)$ moving
$N_0$; normality over $\Q(t_1,\ldots,t_d)$ makes
$\sigma$ preserve $M$.

The field $C_M=M\cap\overline{\Q}$ is finite Galois
over $\Q$, and $\sigma|_{C_M}\ne1$. The normal
finite étale model of F.35 can be chosen
$\sigma$-stable over
\[
 \mathcal O_{C_M}[S^{-1}]
       [t_1,\ldots,t_d,1/f],
       \qquad f\in\Z[t_1,\ldots,t_d],
\]
with the required polynomial derivatives and resultants
units. Choose distinct good rational primes $p_i$
completely split in $C_M$. An unramified Gauss
prolongation $v_i$ over a chosen constant prime
satisfies $v_i(p_i)=1$, and differs from
$\sigma^{-1}v_i$ because their restrictions to
$C_M$ differ. Its residue field is infinite when
$d>0$; when $d=0$, choose $p_i$ larger than the
finite root-avoidance bound.
In the corresponding semilocal Dedekind ring impose
\[
 \alpha\equiv r_i+p_i\pmod{\mathfrak p_{v_i}^2},
 \qquad
 \alpha\equiv c_i\pmod{\mathfrak p_{\sigma^{-1}v_i}},
\]
where $r_i$ is a root of $P_i$ and $c_i$ avoids
every reduced $P_j$-root. All coefficients are fixed
by $\sigma$, although their roots may be permuted.
Taylor expansion and the good resultants give
\[
 v_i(P_j(\alpha))=\delta_{ij},\qquad
 v_i(P_j(\sigma\alpha))=0.
\]
Thus, for $R_j=P_j(\sigma\alpha)/P_j(\alpha)$,
$v_i(R_j)=-\delta_{ij}$, and $\alpha\in U(Mk)$.

Let $b_\alpha\in Mk$ be the actual value at this
point. Choose a finite $S_0/K$ in $k$ such that
$F=MS_0$ contains $b_\alpha$. The field $F$ is
finite over $M$, finitely generated over $\Q$,
$\sigma$-stable, and fixed before any radical
level is chosen.
For $z^n=a\in F^\times$, $z\in Mk$, and
$\tau\in\operatorname{Gal}(Mk/F)$, the two
automorphisms $\sigma,\tau$ commute. On $M$ this
follows because $\tau$ is the identity and
$\sigma$ preserves $M$. On $k$, $\sigma$ is the
identity and $\tau$ preserves $k$: the latter
holds since $\tau$ fixes $K\subset M$ and $k/K$
is Galois. This is the required replacement for
normality over $\Q$ in F.35.
Since $\tau z/z\in\mu_n\subset k$, the earlier
calculation gives $\sigma z/z\in F$, and hence
$\sigma a/a\in(F^\times)^n$.

Apply this to the evaluated relation (F.43.8).
For representatives $a_{j,n}\in\Z$ of
$\lambda_j\bmod n$ one obtains
\[
 \frac{\sigma b_\alpha}{b_\alpha}
       \prod_j R_j^{-a_{j,n}}\in(F^\times)^n .
\]
The integer characters $v_i\circ N_{F/M}$ give
\[
 [F:M]\lambda_i
   =-(v_i\circ N_{F/M})(\sigma b_\alpha/b_\alpha)
       \in\Z\quad\hbox{in }\widehat{\Z}.
\]
It follows that
$\lambda_i\in\Q\cap\widehat{\Z}=\Z$.
Therefore the class is the Kummer class of
\[
 h(T)=b_0\prod_j
       \left(P_j(T)/P_j(x_0)\right)^{\lambda_j}
       \in\mathcal O(U)^\times.                             \tag{F.43.9}
\]
At every finite-extension point its actual value
differs from $h$ by an element of $k^\times$,
by (F.43.5).

\proofparagraph{The finite arithmetic span and generalized Albanese.}
Let $X/L$ be smooth, geometrically connected, affine
and pointed at $x_0$, and let a compatible Kummer
class $\xi$ have actual values at all finite-extension
points. After a finite extension splitting the
boundary of a smooth projective compactification,
descend the data and its generalized Albanese $A$
to an $F$ as in (F.43.6). The integral comparison
used in F.30 is valid in characteristic zero:
\[
 H^1(X_{\bar L},\Z_\ell(1))
       \simeq\operatorname{Hom}(T_\ell A,\Z_\ell(1)).
                                                               \tag{F.43.10}
\]
The finite kernel in the geometric abelian
fundamental-group comparison is killed by Hom
into the torsion-free target
\citep[Propositions~4.1 and~4.3,
Corollary~4.2]{F_SpiessSzamuely2003}.
This gives the integral identification without
a saturation assumption.

Put $H=G_L\triangleleft G_F$ and
$W=(V_\ell A)_H$, with lattice the image of
$T_\ell A$. Since $\mu_\infty\subset L$, the
normalized geometric functional defined by $\xi$
extends to an arithmetic homomorphism on the
split fundamental-group extension, zero on the
section at $x_0$. The projected point Kummer
homomorphisms span $W$. Otherwise a nonzero
annihilating integral functional, reduced
at a finite level where it is nonzero, would
give a nontrivial finite quotient of
$\pi_1(X)$ trivial on every closed-point
decomposition group. This contradicts the
arithmetic finite-quotient detection above.

Choose finitely many such point images spanning
$W$. Their coordinates and actual values
$\beta_i$, as well as the normalizing value
$b_0$, lie in a finitely generated
$M/\Q$ finite over a suitable descent field,
by (F.43.6). Restriction to $G_{Mk}$ preserves
their rational spans: this group has finite
index in each selected point group, and the
projected Kummer maps are homomorphisms.
For $g\in G_M$ and $h\in G_{Mk}$,
\[
 c_i(ghg^{-1})=g c_i(h),\qquad
 \kappa(\beta_i/b_0)(ghg^{-1})
       =\chi(g)\kappa(\beta_i/b_0)(h).
\]
The points and values are defined over $M$;
the first identity's coboundary term vanishes
in the $H$-coinvariants. Normality of $G_L$
in $G_F$ preserves those coinvariants.
Equality with the prescribed evaluations
makes the functional $G_M$-equivariant on
the spanning images, hence on all of $W$.

The integral Tate-Hom theorem for $1$-motives
over the finitely generated field $M$ applies:
the exact original is
\citet[Theorem~4.3(ii), p.~676]{F_Jannsen1995}.
Apply it to $[0\to A]$ and $[0\to\mathbb G_m]$.
It places this functional in
$\operatorname{Hom}_M(A,\mathbb G_m)\otimes\Z_\ell$.
After the boundary is split, its normalized
geometric unit lattice is fixed over the
descent field: the boundary divisor and
value $1$ at $x_0$ determine a unit uniquely.
The Albanese universal property therefore
puts $\xi$ in this one unit lattice tensored
with $\widehat{\Z}$, plus the actual constant
value at $x_0$.
This eliminates the Picard--Tate component
of the given class; it does not assert that
the full Picard--Tate module is zero.

For a curve, F.27's Riemann--Roch finite
map to a punctured rational line isolates
each boundary residue. Its norm preserves
actual point values, and (F.43.9) makes
the norm residue an ordinary integer.
The residue is weighted by residue degree,
not ramification index; the selected
rational boundary point has degree one.
The injective integral divisor matrix
of a unit basis then makes every unit
exponent an integer.
In higher dimension choose smooth transverse
complete-intersection curves through smooth
points of the boundary components and
away from their intersections. At a finite
Kummer level, after étale localization,
the class is represented by $u z^a$.
Transversality carries $z$ to a uniformizer
and preserves its residue. The curve
conclusion therefore makes all boundary
residues integral, and the same divisor
matrix gives integral exponents.

Normalize the resulting regular unit $h$
by $h(x_0)=b_0$. A finite Galois extension
used above introduces no descent ambiguity:
the quotient of a conjugate of $h$ by $h$
has zero divisor on a proper completion,
so is constant, and its value at $x_0$ is
one. Thus $h$ descends. At every finite
Kummer level the residual class is
geometrically zero; the low-degree
Hochschild--Serre sequence makes it
constant, and evaluation at $x_0$ kills it.
This proves $\xi=\kappa(h)$ over $L$.

\proofparagraph{Point functions and the precise CHW hypothesis.}
For a pinned relative map from a smooth
affine curve to $\mathbb G_m$, we have
\[
                  \phi(x)=c_xh(x),\qquad c_x\in k^\times.
                                                               \tag{F.43.11}
\]
These scalars may now be transcendental
over $\Q$. The pole-degree argument of
F.28 uses only that they belong to the
fixed subfield $k\subset L$.
If $h$ is nonconstant, points of degree
larger than $\deg(h)$ show that the
point function cannot be constant.
Pull back $T-1$ on the nonempty
inverse-image open, obtaining an actual
unit $h_1$. Choose there a closed point
whose degree exceeds the combined pole
bound for $h,h_1$. The function
$c_xh-d_xh_1-1$ vanishes at that point,
so it vanishes identically. Consequently
$h_1=Ah+B$ with $A,B\in k^\times$.
At every point where $h(x)\notin k$,
the two coefficients in the resulting
linear identity give
$c_x=-A/B=:c$.
For every $b\in k^\times$, repeat the
argument on the inverse-image open of
$T-b\ne0$ to obtain
$h_b=e_b(ch-b)$.
If $h(x)\in k$, choose $b=ch(x)\ne0$.
Were $\phi(x)\ne b$, this point would
belong to that open while its regular
unit $h_b$ vanished. Thus $\phi=ch$
at these points as well.

If $h$ is constant, (F.43.11) makes
all values lie in $hk^\times\subset L$.
By (F.43.3), after a finite extension
there is $\beta$ with $\beta/h\notin k$.
All preceding results persist over that
extension. No closed point maps to
$\beta$, so the whole curve factors
through the complement of $\beta$:
the inverse image of that closed point
is closed, and every nonempty closed
subset of a noetherian curve has a
closed point. Pulling back $T-\beta$
therefore gives an actual unit
$h_\beta$ on the entire curve.
Its values at closed points belong to
the base field, since its actual
evaluations and residual scalars do.
The unbounded closed-point degrees
force $h_\beta$ to be constant.
The value ratios consequently give
\[
              \frac{\phi(x)-\beta}
                   {\phi(x_0)-\beta}\in k.
\]
Together with (F.43.11) and
$\beta/h\notin k$, this forces
$\phi(x)=\phi(x_0)$.

For a higher-dimensional $X$, make
an arbitrary geometric point rational
by a finite extension, and choose a
smooth geometrically connected curve
through it and $x_0$. The curve's
point-function unit and $h$ restricted
to the curve have the same Kummer
class. Their quotient has zero divisor
on a proper completion, hence is
constant, and its value at $x_0$ is
one. Thus they agree. The entire
geometric-point function is the unique
regular unit normalized at $x_0$;
uniqueness supplies finite descent.
The dimension-zero case is the
field-source statement.

This meets exactly
\citet[Theorem~4.18(5)]{F_CHW2024}:
a connected regular affine source
with a rational point and a pinned
relative map to $\mathbb G_m$ has
an algebraic geometric-point function.
In characteristic zero regularity
gives smoothness, and connectedness
together with a rational point gives
geometric connectedness. That theorem
has no additional amenability
hypothesis. It supplies the full
étale-chart and descent comparison,
and hence the all-scheme reconstruction
statement over every finitely generated
$L/k$. Equality on points alone is
not substituted for that comparison.

\proofparagraph{Arbitrary subfields by relative constants.}
Now let $k_0\subset L$ be any subfield.
Choose a transcendence basis $B$ of
$k_0/\Q$; it is finite because
$L/\Q$ has finite transcendence degree.
Put $F_0=\Q(B)\subset k_0$.
Choose a finitely generated $F\subset L$
containing $K,B$ and generators for
$L/k$, so that $L=Fk$ is Galois
prosolvable over $F$. Let $C$ be the
relative algebraic closure of $F_0$
in $F$. It is finite over $F_0$ and
finitely generated over $\Q$, and
$F/C$ is regular.

For $a\in k_0$, let $N/C$ be the
normal closure of $C(a)/C$. It is
finite since $a$ is algebraic over
$F_0$. The defining property of $C$
gives $N\cap F=C$. Equation (F.43.1)
shows that the conjugates of $a$
over $C$ are its conjugates over
$F$. Normality of $L/F$ puts
$N$ inside $L$, and
\[
 \operatorname{Gal}(N/C)
       \simeq\operatorname{Gal}(FN/F)
\]
is a finite quotient of the prosolvable
group $\operatorname{Gal}(L/F)$.
It is solvable. Hence every $a\in k_0$
lies in $C^{\mathrm{solv}}$, giving
\[
                    k_0\subset C^{\mathrm{solv}}.           \tag{F.43.12}
\]
This extension is algebraic: both
fields are algebraic over
$F_0\subset k_0$.
The core theorem just proved for
the finitely generated field $C/\Q$
gives reconstruction over
$C^{\mathrm{solv}}$.
Algebraic-extension descent F.32
now proves reconstruction over
$k_0$. The possibly transcendental
extension $k_0k/k_0$ was not used.

\proofparagraph{The arithmetic description and strict enlargement.}
Suppose first that the normal closure
$N/F_0$ of an algebraic
$k_0/F_0$ has an open prosolvable
subgroup $H$. Its fixed field
$C=N^H$ is finite over $F_0$,
and $N/C$ is Galois prosolvable.
Thus $k_0\subset N\subset
C^{\mathrm{solv}}$, which gives
both the asserted embedding class
and reconstruction by F.32.

Conversely, (F.43.12) places any
field in the embedding class in
$C^{\mathrm{solv}}$, with
$C/F_0$ finite and
$F_0\subset k_0$ finitely
generated over $\Q$. Let
$C'/F_0$ be the finite normal
closure of $C$. There are only
finitely many $F_0$-conjugates
of $C^{\mathrm{solv}}$, since
it is normal over $C$.
Their compositum $E/F_0$ is
Galois and contains $C'$.
After base change to $C'$ each
conjugate is Galois prosolvable.
The group of their finite
compositum embeds as a closed
subgroup of the product of
these prosolvable groups, so
$E/C'$ is prosolvable.
Consequently $\operatorname{Gal}(E/F_0)$
has an open prosolvable subgroup.
The normal closure of $k_0/F_0$
lies in $E$, and its group is
a quotient of this virtually
prosolvable group. It too is
virtually prosolvable, proving
the equivalence.

Finally take $K=\Q(t)$.
If $\Q(t)^{\mathrm{solv}}$
embedded in a finitely generated
$E/\Q^{\mathrm{solv}}$, divisibility
would put the image of every
nonzero element, including $t$,
in $D(E^\times)$.
F.33 identifies the latter with
$(\Q^{\mathrm{solv}})^\times$,
whose elements are algebraic
over $\Q$. A characteristic-zero
field embedding fixes $\Q$ and
preserves transcendence, a
contradiction. This proves the
strict enlargement of F.37.
The intrinsic converse problem
of F.9 remains unresolved.
\end{proof}

The targeted literature search recorded in
\path{sources/F_sources.md} found no opened primary source
stating the sufficient class of F.43. The current
Carlson--Stix theorem still assumes a sub-$p$-adic
field \citep[Theorem~A]{F_CS2026}, and the supplied
CHW theorem is used here through its general reduction
criterion. This is a comparison of the checked
statements, not a bibliographic priority claim.
